\subsection{Evaluation}
To evaluate the outcomes of the simulations and compare the success of each policy regime, four key performance metrics were defined, reflecting the multi-faceted nature of effective local governance.

\begin{enumerate}
    \item \textbf{Maximum Sustainable Compliance:} Defined as the highest population-weighted average compliance rate the LGU agent managed to achieve and maintain over the long term~\citep{ApostolJamoralin2024, Torkayesh2021}. This established the agent's ability to solve the core environmental problem over a sustained period without experiencing behavioral collapse.
    \item \textbf{Cost-Effectiveness:} Calculated as the average percentage of compliance gained per \textpeso~100,000 spent. By integrating this financial ratio, the metric addressed the need for fiscal discipline and efficiency, ensuring the policy provided the maximum impact for its cost---a major practical concern for budget-constrained LGUs~\citep{MedinaMijangos2020, WorldBank2022}.
    \item \textbf{Policy Equity:} Assessed by measuring the final variance in compliance rates between the 7 barangays~\citep{Jimenez2025}. A lower variance indicated a more equitable and evenly distributed policy outcome, verifying that the DRL agent's resource allocation strategy did not simply concentrate resources on easy-to-manage areas, but rather succeeded in raising compliance across all geographic and socioeconomic units~\citep{TianReview2024}.
    \item \textbf{Optimal Resource Allocation:} Analyzed for the comprehensive Hybrid regime as a key prescriptive output. This metric detailed the final learned budget split, in both percentage (\%) and peso (\textpeso) terms, across the three policy levers (IEC, Enforcement, Incentives) and all 7 barangays. This provided essential, actionable intelligence, revealing the data-driven policy mix deemed optimal for the LGU's specific context~\citep{MuZhang2021, Zhao2022}.
\end{enumerate}

\subsection{Model Validation and Verification}
\setlength{\parindent}{2em}

The validation of the BacolodGymEnv follows a structural and face validity protocol rather than a predictive or empirical validity protocol. This distinction is critical given that longitudinal, high-frequency data on household-level waste compliance in the municipality is currently unavailable. Consistent with the strict operational optimization framing of this study, we explicitly do not claim predictive validity. Instead, the integrated model is validated purely as a controlled computational testbed designed for comparative policy analysis.

The model is verified and validated through three primary mechanisms: face validity, structural validity, and replicability. Face validity is established by ensuring that the model's initial conditions—such as municipal and barangay budget limits, population densities, and baseline compliance estimates—are strictly calibrated to match official municipal records, qualitative key informant estimates, and reported figures from the Bacolod LGU 10-year Solid Waste Management Plan. Structural validity is confirmed through Global Sensitivity Analysis (as detailed in Section 4.2), which verifies that the model's internal causal pathways and mechanisms behave consistently with established behavioral economics theories. Specifically, the analysis confirms expected patterns, such as the computational "Cost of Effort" negatively affecting compliance, and financial incentives exerting a positive influence subject to diminishing returns. Furthermore, replicability is ensured through the comprehensive documentation of all simulation code and parameters, with all reported outcomes representing averages across multiple random seeds to mitigate stochastic variance.

Acknowledging these operational boundaries, it is crucial to delineate what is not validated within this framework. The model has not been tested against out-of-sample, real-world compliance time series, nor does it attempt to provide precise time-series forecasting. Consequently, no claim is made that the highly optimized simulated compliance rates achieved by the HuDRL agent would be perfectly mirrored in reality. The simulation serves purely as a computational laboratory. Its fundamental value lies in identifying the structural strengths and weaknesses of different policy regimes, establishing a relative ordering of policy effectiveness, and uncovering the underlying mathematical logic of the discovered "Sequential Saturation" strategy, all of which provide a rigorously tested foundation to inform real-world pilot designs.
\subsection{Structural Robustness and Sensitivty Evaluation}
To validate the structural integrity of the proposed HuDRL-driven policy, a hardcoded robustness test was designed to evaluate the model's sensitivity to variations in social tipping points. While the behavioral thresholds for the ``Social Norm Shield'' were initially calibrated based on established social complexity theories~\cite{Centola2018}, it is critical to determine if the policy's success is contingent upon specific quantitative calibrations of these parameters.

The robustness test utilizes a 9-scenario matrix that systematically varies the two primary behavioral thresholds: the compliance floor ($T_{low}$) and the social norm tipping point ($T_{high}$). These variables represent the degree of social inertia and the threshold required to trigger the ``Habit Shield'' within a barangay, respectively. The testing matrix is defined as:

\begin{table}[htbp]
\centering
\caption{Robustness Testing Matrix ($\mathcal{M}$) Parameters}
\label{tab:sensitivity_matrix}
\begin{tabular}{@{}cccc@{}}
\toprule
\textbf{Scenario} & \textbf{\(T_{low}\) (Inertia)} & \textbf{\(T_{high}\) (Shield)} & \textbf{Environment Type} \\ \midrule
1                 & 0.20                           & 0.60                            & Highly Optimistic         \\
2                 & 0.25                           & 0.70                            & Baseline Optimistic      \\
3                 & 0.30                           & 0.80                            & Stubborn (High Friction)  \\
4                 & 0.35                           & 0.60                            & Moderately Optimistic    \\
5                 & 0.40                           & 0.70                            & Standard                 \\
6                 & 0.45                           & 0.80                            & Stubborn (Medium Inertia)\\
7                 & 0.50                           & 0.70                            & High Inertia             \\
8                 & 0.55                           & 0.80                            & Highly Stubborn          \\
9                 & 0.60                           & 0.85                            & Extreme Stubbornness      \\ \bottomrule
\end{tabular}
\end{table}

The scenarios were selected to cover a spectrum of behavioral environments:
\begin{itemize}
    \item \textbf{Optimistic Scenarios:} Lower \(T_{high}\) values (e.g., 0.60) where social norms are easily shifted.
    \item \textbf{Stubborn Scenarios:} Higher \(T_{high}\) values (e.g., 0.80--0.85) representing communities with high social friction and resistance to policy change.
\end{itemize}

\subsection{Computational Setup and Multiprocessing}
Given the stochastic nature of the Agent-Based Model (ABM), each of the nine scenarios was executed for a simulated duration of 1,080 days (12 quarters). To ensure computational efficiency and statistical independence, the robustness test was implemented using a Python-based multiprocessing framework~\cite{mesa2023}. 

Each scenario was assigned to an independent CPU core, where a unique instance of the \texttt{BacolodModel} was initialized with a specific \((T_{low}, T_{high})\) pair injected into the \texttt{HouseholdAgent} behavioral parameters. The HuDRL policy was maintained as the active governance logic across all runs to isolate the impact of behavioral variance on the agent's performance. The final municipality-wide compliance rate was then recorded for comparative analysis in Chapter~4.